% Options for packages loaded elsewhere
\PassOptionsToPackage{unicode}{hyperref}
\PassOptionsToPackage{hyphens}{url}
\PassOptionsToPackage{dvipsnames,svgnames,x11names}{xcolor}
\documentclass[
  10pt,
]{article}
\usepackage{xcolor}
\usepackage[margin=2.4cm]{geometry}
\usepackage{amsmath,amssymb}
\setcounter{secnumdepth}{-\maxdimen} % remove section numbering
\usepackage{iftex}
\ifPDFTeX
  \usepackage[T1]{fontenc}
  \usepackage[utf8]{inputenc}
  \usepackage{textcomp} % provide euro and other symbols
\else % if luatex or xetex
  \usepackage{unicode-math} % this also loads fontspec
  \defaultfontfeatures{Scale=MatchLowercase}
  \defaultfontfeatures[\rmfamily]{Ligatures=TeX,Scale=1}
\fi
\usepackage{lmodern}
\ifPDFTeX\else
  % xetex/luatex font selection
  \setmainfont[]{Libertinus Serif}
  \setmonofont[Scale=0.85]{Latin Modern Mono}
  \setmathfont[]{Latin Modern Math}
\fi
% Use upquote if available, for straight quotes in verbatim environments
\IfFileExists{upquote.sty}{\usepackage{upquote}}{}
\IfFileExists{microtype.sty}{% use microtype if available
  \usepackage[]{microtype}
  \UseMicrotypeSet[protrusion]{basicmath} % disable protrusion for tt fonts
}{}
\makeatletter
\@ifundefined{KOMAClassName}{% if non-KOMA class
  \IfFileExists{parskip.sty}{%
    \usepackage{parskip}
  }{% else
    \setlength{\parindent}{0pt}
    \setlength{\parskip}{6pt plus 2pt minus 1pt}}
}{% if KOMA class
  \KOMAoptions{parskip=half}}
\makeatother
\usepackage{longtable,booktabs,array}
\newcounter{none} % for unnumbered tables
\usepackage{calc} % for calculating minipage widths
% Correct order of tables after \paragraph or \subparagraph
\usepackage{etoolbox}
\makeatletter
\patchcmd\longtable{\par}{\if@noskipsec\mbox{}\fi\par}{}{}
\makeatother
% Allow footnotes in longtable head/foot
\IfFileExists{footnotehyper.sty}{\usepackage{footnotehyper}}{\usepackage{footnote}}
\makesavenoteenv{longtable}
\setlength{\emergencystretch}{3em} % prevent overfull lines
\providecommand{\tightlist}{%
  \setlength{\itemsep}{0pt}\setlength{\parskip}{0pt}}

\providecommand{\E}{\mathbb{E}}
\providecommand{\N}{\mathbb{N}}
\providecommand{\Z}{\mathbb{Z}}
\providecommand{\R}{\mathbb{R}}
\providecommand{\eps}{\varepsilon}
\providecommand{\ceil}[1]{\left\lceil #1\right\rceil}
\providecommand{\floor}[1]{\left\lfloor #1\right\rfloor}
\AtBeginDocument{\let\setminus\smallsetminus}
\usepackage[most]{tcolorbox}
\newtcolorbox{warning}{colback=red!5,colframe=red!65!black,boxrule=0.8pt,arc=2pt,left=6pt,right=6pt,top=5pt,bottom=5pt}
\usepackage{etoolbox}
\AtBeginEnvironment{longtable}{\small}
\setlength{\emergencystretch}{3em}
\usepackage{fancyhdr}
\pagestyle{fancy}
\fancyhf{}
\fancyhead[L]{\small\bfseries UNREFEREED GPT-6 ASTRA OUTPUT}
\fancyhead[R]{\small\thepage}
\renewcommand{\headrulewidth}{0.4pt}
\usepackage{bookmark}
\IfFileExists{xurl.sty}{\usepackage{xurl}}{} % add URL line breaks if available
\urlstyle{same}
\hypersetup{
  pdftitle={Problem 5.1},
  pdfauthor={Writeup: gpt-6-astra (single model pass)},
  colorlinks=true,
  linkcolor={blue!50!black},
  filecolor={Maroon},
  citecolor={Blue},
  urlcolor={blue!50!black},
  pdfcreator={LaTeX via pandoc}}

\title{Problem 5.1}
\usepackage{etoolbox}
\makeatletter
\providecommand{\subtitle}[1]{% add subtitle to \maketitle
  \apptocmd{\@title}{\par {\large #1 \par}}{}{}
}
\makeatother
\subtitle{Unrefereed candidate proof by GPT-6 Astra}
\author{Writeup: \texttt{gpt-6-astra} (single model pass)}
\date{Catalog id \texttt{2507.10840\_\_01} --- generated 2026-09-17}

\begin{document}
\maketitle

\begin{warning}

\textbf{UNREFEREED MODEL OUTPUT.} This document was selected solely
because GPT-6 Astra labelled its own result
\texttt{would\_publish:\ true} and returned \texttt{proved} or
\texttt{disproved}. It has not passed the adversarial LLM referee used
for the earlier Sol campaign, has not been checked by a human
mathematician, and has not been checked for novelty. Treat every
mathematical and bibliographic claim below as unverified.

\end{warning}

\section{Summary and provenance}\label{summary-and-provenance}

{\def\LTcaptype{none} % do not increment counter
\begin{longtable}[]{@{}
  >{\raggedright\arraybackslash}p{(\linewidth - 2\tabcolsep) * \real{0.5000}}
  >{\raggedright\arraybackslash}p{(\linewidth - 2\tabcolsep) * \real{0.5000}}@{}}
\toprule\noalign{}
\endhead
\bottomrule\noalign{}
\endlastfoot
Catalog id & \texttt{2507.10840\_\_01} \\
Catalog entry &
\href{https://graph-theory-ai.github.io/graph-conjectures/arxiv/2507.10840__01/}{Problem
5.1} \\
Source corpus & arxiv \\
Campaign leg & selected second attempts (\texttt{attacks\_retry}) \\
Source paper / entry & Covering Complete Geometric Graphs by Monotone
Paths \\
Model verdict & \textbf{proved} (confidence: high) \\
Model's one-line claim & An odd regular polygon with a central cluster
yields point sets requiring at least 3n/4 - O(sqrt(n)) crossing-free
paths. \\
Model & \texttt{gpt-6-astra}, reasoning effort \texttt{max},
\texttt{mode=pro}, flex service tier \\
Original artifact &
\texttt{attacks\_retry/2507.10840\_\_01/output.md} \\
Independent review & \textbf{None} \\
Model's caveats & The optimal leading constant is not determined; paths
have their standard graph-theoretic meaning. \\
\end{longtable}
}

\section{Problem statement}\label{problem-statement}

\emph{No extracted statement was stored in the artifact metadata.}

\section{Astra writeup}\label{astra-writeup}

\emph{The text below is the model output verbatim, apart from moving its
machine-readable verdict block into the summary above and shifting
Markdown heading levels for this document. Mechanical TeX defects and
equation tags were normalized where needed for compilation.}

\subsection{An affirmative
construction}\label{an-affirmative-construction}

Write \(\pi(A)\) for the minimum number of crossing-free paths covering
all edges of \(K_n[A]\). Covering paths may overlap, and they need not
be Hamiltonian.

We prove the following.

\textbf{Theorem.} 1. For every positive integer \(q\), there is a
\(6q\)-point set \(A\) in general position such that \[
   \pi(A)\ge 4q=\frac23|A|.
   \] 2. More strongly, for every sufficiently large \(n\), there is an
\(n\)-point set \(A\) in general position such that \[
   \pi(A)\ge \frac34n-O(\sqrt n).
   \]

Thus \(c=2/3\) answers the question affirmatively. In fact, every fixed
\(c<3/4\) works for all sufficiently large \(n\).

The key is a simultaneous incompatibility between long chords of an odd
polygon and edges incident with points near its center. The geometric
and counting claims needed below are proved directly.

\subsubsection{1. The geometric gadget}\label{the-geometric-gadget}

Let \[
h=2m+1\ge 5
\] be odd, and let \[
v_j=\left(\cos\frac{2\pi j}{h},\ \sin\frac{2\pi j}{h}\right),
\qquad 0\le j<h.
\] Indices are taken modulo \(h\). Put \(V=\{v_0,\ldots,v_{h-1}\}\), and
define the following set of \(h\) chords: \[
D=\{v_i v_{i+m}:0\le i<h\}.
\]

We need two properties.

\paragraph{\texorpdfstring{Property 1: disjoint edges of \(D\)
cross}{Property 1: disjoint edges of D cross}}\label{property-1-disjoint-edges-of-d-cross}

Fix \(v_0v_m\). Any edge of \(D\) disjoint from it has one endpoint in
\[
\{v_1,\ldots,v_{m-1}\}
\] and the other in \[
\{v_{m+1},\ldots,v_{2m}\}.
\] This follows immediately by checking the indices in \(v_i v_{i+m}\).
Its endpoints therefore alternate with \(v_0,v_m\) around the convex
polygon, so the two chords cross. Rotation gives the assertion for every
edge of \(D\).

Consequently, every crossing-free path contains at most two edges of
\(D\): any three edges of a simple path contain two vertex-disjoint
edges. If a crossing-free path contains two edges of \(D\), they must
share a vertex.

The two edges of \(D\) incident with \(v_i\) are precisely \[
v_i v_{i+m}
\quad\text{and}\quad
v_i v_{i+m+1}.
\]

\paragraph{Property 2: two incident chords trap a central
region}\label{property-2-two-incident-chords-trap-a-central-region}

Define \[
T_i=\operatorname{conv}\{v_i,v_{i+m},v_{i+m+1}\}.
\] The origin lies strictly inside every \(T_i\). Indeed, after rotating
\(v_i\) to \((1,0)\), the other two vertices have coordinates \[
\left(-\cos\frac{\pi}{h},\ \sin\frac{\pi}{h}\right),
\qquad
\left(-\cos\frac{\pi}{h},\ -\sin\frac{\pi}{h}\right).
\] Hence \[
U=\bigcap_{i=0}^{h-1}\operatorname{int}(T_i)
\] is an open neighborhood of the origin.

Suppose that a crossing-free path contains \[
v_i a,\ v_i b,
\qquad a=v_{i+m},\quad b=v_{i+m+1}.
\] For any \(x\in U\) and any \[
z\in V\setminus\{v_i,a,b\},
\] the segment \(xz\) crosses one of these two chords.

To see this, \(x\) lies inside \(T_i\), whereas \(z\), being another
vertex of a strictly convex polygon, lies outside \(T_i\). The segment
\(xz\) must leave \(T_i\). It cannot cross \(ab\), since \(ab\) is a
side of the outer convex polygon and both \(x\) and \(z\) lie strictly
on its inner side. It must therefore cross \(v_i a\) or \(v_i b\). With
the points in general position, this is a proper crossing.

Thus, when both chords are present, edges from central points to outer
vertices are possible only at the three vertices \(v_i,a,b\).

\subsubsection{\texorpdfstring{2. A particularly simple proof with
\(c=2/3\)}{2. A particularly simple proof with c=2/3}}\label{a-particularly-simple-proof-with-c23}

Take \[
A=V\cup\{o\},
\qquad o=(0,0),
\] and let \[
S=\{ov_j:0\le j<h\}
\] be the central star.

This point set is in general position. Three points on the circle cannot
be collinear, and a line through the origin contains two of the polygon
vertices only if they are antipodal, which is impossible because \(h\)
is odd.

We claim that every crossing-free path \(P\) satisfies \[
|E(P)\cap(D\cup S)|\le 3. \qquad\text{(1)}
\]

If \(P\) contains at most one edge of \(D\), this follows from \[
\deg_P(o)\le 2.
\]

Otherwise, its two edges of \(D\) are \(v_i a,v_i b\), as above.
Property 2 says that a spoke in \(P\) can only be one of \[
ov_i,\quad oa,\quad ob.
\] The edge \(ov_i\) is impossible because \(v_i\) already has degree
two in \(P\). Moreover, \(oa\) and \(ob\) cannot both occur, since
together with \(v_i a,v_i b\) they form the cycle \[
o,a,v_i,b,o.
\] A simple path cannot contain a cycle. Hence \(P\) contains at most
one spoke, proving (1).

Now \(D\) and \(S\) are disjoint and each has \(h\) edges. If \(k\)
crossing-free paths cover \(K_{h+1}[A]\), then, even allowing overlaps,
\[
3k\ge |D\cup S|=2h.
\] Therefore \[
\pi(A)\ge \left\lceil\frac{2h}{3}\right\rceil. \qquad\text{(2)}
\]

Choose \(h=6q-1\). Then \(n=h+1=6q\), and (2) gives \[
\pi(A)\ge
\left\lceil\frac{12q-2}{3}\right\rceil
=4q
=\frac23n.
\] This already proves the requested assertion.

\subsubsection{\texorpdfstring{3. A central cluster improves the
coefficient to
\(3/4-o(1)\)}{3. A central cluster improves the coefficient to 3/4-o(1)}}\label{a-central-cluster-improves-the-coefficient-to-34-o1}

We now replace the single center by \(r\ge2\) points.

Choose a set \(B\subset U\) of \(r\) points such that \[
A=V\cup B
\] is in general position. This is possible by choosing the points
successively in the open set \(U\), avoiding the finitely many lines
determined by previously chosen points.

Let \[
S=\{vx:v\in V,\ x\in B\},
\] so that \[
|S|=hr.
\]

For a crossing-free path \(P\), write \[
d(P)=|E(P)\cap D|,
\qquad
s(P)=|E(P)\cap S|.
\] We have \[
d(P)\le2,\qquad s(P)\le2r. \qquad\text{(3)}
\] The second inequality follows by summing the path degrees at the
\(r\) central vertices.

There is a much stronger bound when \(d(P)=2\): \[
d(P)=2\quad\Longrightarrow\quad s(P)\le2. \qquad\text{(4)}
\] Indeed, let the two chords be \(v_i a,v_i b\). Property 2 restricts
all spokes in \(P\) to outer endpoints in \(\{v_i,a,b\}\). The vertex
\(v_i\) already has degree two, while \(a\) and \(b\) each already have
degree one. Thus at most one spoke can use \(a\), and at most one can
use \(b\).

Assign weight \(2r-2\) to each edge of \(D\), weight \(1\) to each edge
of \(S\), and weight \(0\) to every other edge. All other edges remain
allowed in the paths.

By (3)--(4), every crossing-free path has weight at most \(4r-2\).
Explicitly, \[
(2r-2)d(P)+s(P)\le
\begin{cases}
2r,&d(P)=0,\\
(2r-2)+2r=4r-2,&d(P)=1,\\
2(2r-2)+2=4r-2,&d(P)=2.
\end{cases}
\] The total edge weight is \[
(2r-2)|D|+|S|
=(2r-2)h+hr
=(3r-2)h.
\] Summing weights over any covering family therefore proves \[
\boxed{\displaystyle
\pi(A)\ge
\left\lceil\frac{(3r-2)h}{4r-2}\right\rceil.}
\qquad\text{(5)}
\]

Finally, for each sufficiently large \(n\), choose an integer \[
r=\sqrt n+O(1)
\] such that \(h=n-r\) is odd. Then \(r\ge2\) and \(h\ge5\), so the
construction applies. Since \[
\frac{3r-2}{4r-2}
=\frac34-\frac{1}{8r-4},
\] equation (5) yields \begin{equation*}
\begin{aligned}
\pi(A)
&\ge
\left(\frac34-\frac{1}{8r-4}\right)(n-r)\\
&=
\frac34n-\frac34r-\frac{n-r}{8r-4}\\
&=
\frac34n-O(\sqrt n).
\end{aligned}
\end{equation*}

This proves the stronger assertion. The argument covers arbitrary
crossing-free paths, including short paths, and allows arbitrary overlap
between covering paths. No assumption of a Hamilton-path decomposition
is needed.

\end{document}
